Mi opinion es que el trabajo ya tiene una base tecnica suficientemente buena como para justificar una lectura cientifica exploratoria seria. La separacion entre valores observados, derivados e imputados es una fortaleza real del reporte, porque evita vender la topologia como si toda proviniera de mediciones directas. Tambien me parece acertado que el documento no intente forzar una ``clasificacion final'' de exoplanetas; en su estado actual, Mapper esta funcionando mejor como instrumento para priorizar regiones y grafos que como veredicto taxonomico.

Si tuviera que resumir el valor del trabajo en una frase, diria que el espacio orbital parece contener la senal mas prometedora y mas honesta del analisis. No es el grafo mas interesante solo porque tenga muchos ciclos, sino porque esos ciclos aparecen donde la imputacion pesa poco. Eso le da una ventaja fuerte frente al espacio termico, cuya complejidad visual es atractiva pero todavia demasiado dependiente de datos completados.

Tambien creo que el resultado sobre \texttt{pl\_dens} es mas importante de lo que parece a primera vista. Que la densidad reduzca algo de la complejidad en lugar de aumentarla es una pista de que el pipeline esta detectando redundancia fisica real y no solo acumulando variables. Ese tipo de sensibilidad es valiosa porque demuestra que Mapper no esta respondiendo de manera ciega a cualquier nueva coordenada.

Mi principal critica constructiva es que el reporte ya merece una capa mas fuerte de interpretacion nodo por nodo. Las figuras muestran que existen regiones centrales, periferias y pequenos componentes aislados, pero el argumento cientifico todavia no termina de cerrar mientras no se traduzcan esas regiones en descripciones fisicas mas concretas y auditables. En otras palabras, el trabajo ya paso la etapa de ``pipeline que corre'' y esta listo para entrar en la etapa de ``estructuras que podemos defender''.
